% Modelo de Avaliação em LaTeX
% IFPR - Paranavaí
% Se utilizar o Overleaf, clique em RICH TEXT acima para uma visualização mais suave do código!

\documentclass[a4paper,11pt]{article}

%-------------------------------------------------------------------------
% PACOTES (não é necessário alterar nada) 
%-------------------------------------------------------------------------

\usepackage[utf8]{inputenc}
\usepackage[brazil]{babel} 
\usepackage{graphicx}
\usepackage{color} 
\usepackage{hyperref} 
\usepackage{geometry} 
\usepackage{amsmath,amsthm,amsfonts,amssymb,amscd,amsxtra} 
\usepackage{multicol}
\usepackage{multirow}
\usepackage{enumerate}
\usepackage{tikz}
\usetikzlibrary{positioning,shapes,fit,arrows,through,calc,shapes.geometric,patterns,decorations.markings}
\usepackage{pgf,pgfplots}
\usepgfplotslibrary{fillbetween}
\pgfplotsset{compat=1.15}
\usepackage{mathrsfs}
\definecolor{myblue}{RGB}{56,94,141}
\definecolor{mygreen}{RGB}{51,153,0}
\definecolor{myred}{RGB}{204,0,0}
\usepackage{setspace}
\geometry{a4paper,left=1.5cm,right=1.5cm,top=1.5cm,bottom=1.5cm}

%-------------------------------------------------------------------------
% INFORMAÇÕES DA AVALIAÇÃO (PREENCHA)
%-------------------------------------------------------------------------

\newcommand{\professor}{Seu nome aqui}
\newcommand{\disciplina}{Nome da Disciplina}
\newcommand{\trimestre}{1º}
\newcommand{\curso}{Nome do Curso}
\newcommand{\turno}{Matutino}
\newcommand{\instrumento}{Tipo de Instrumento Avaliativo} % Pode ser Prova Escrita, Trabalho em classe, Trabalho Extraclasse, Relatório, etc.

%-------------------------------------------------------------------------
% CABEÇALHO (não é necessário alterar nada) 
%-------------------------------------------------------------------------

\begin{document}

\pagestyle{empty}

\begin{minipage}{6cm} %Logo do IFPR
\resizebox{6cm}{1.6cm}{
\begin{tikzpicture}[square/.style={regular polygon,regular polygon sides=4}]
\tikzset{mynode/.style={square,rounded corners,draw=mygreen, top color=mygreen, bottom color=mygreen,very thick, text centered},}  
\node[circle, draw=myred , top color=myred, bottom color=myred] (q1) {$\quad$};
\node[mynode, below=0.08cm of q1] (q2) {$\;\;$};
\node[mynode, below=0.08cm of q2] (q3) {$\;\;$};
\node[mynode, below=0.08cm of q3] (q4) {$\;\;$};
\node[mynode, right=0.08cm of q1] (q5) {$\;\;$};
\node[mynode, below=0.08cm of q5] (q6) {$\;\;$};
\node[mynode, below=0.08cm of q6] (q7) {$\;\;$};
\node[mynode, below=0.08cm of q7] (q8) {$\;\;$};
\node[mynode, right=0.08cm of q5] (q9) {$\;\;$};
\node[mynode, right=0.08cm of q7] (q10) {$\;\;$};
\node[right=0.08cm of q8] (aux) {};
\node[text centered, right=0.06cm of q10] (q11) {\huge \sf \textbf{INSTITUTO FEDERAL}};
\node[mygreen, text centered, right=0.5cm of aux] (q12) {\huge\textsf{Paraná}};
\end{tikzpicture}}
\end{minipage}
\begin{minipage}{7cm}
\begin{center}
\begin{footnotesize}
\sf \textbf{INSTITUTO FEDERAL DO PARANÁ}\\
\textbf{Campus Paranavaí}\\
Rua José Felipe Tequinha, 1400\\
Jardim das Nações - Paranavaí - PR
\end{footnotesize}
\end{center}
\end{minipage}
\begin{minipage}{4cm}
\centering
\includegraphics[scale=0.12]{meBrasil.png}\\
\begin{scriptsize}
\textsf{Ministério da Educação}
\end{scriptsize}
\end{minipage}
\doublespacing
\begin{center}
\begin{tabular}{|p{6.5cm}|p{6.5cm}|p{3.5cm}|}
\hline
\textbf{Curso:} \curso & \textbf{Turno:} \turno & \textbf{Trimestre:} \trimestre \\
\hline
\textbf{Disciplina:} \disciplina & \textbf{Professor(a):} \professor & \textbf{Data:} \\
\hline
\multicolumn{2}{|l|}{\textbf{Nome:}} & \textbf{R.A.:}\\
\hline
\textbf{Conceito:} $\square$ A $\;\;\square$ B $\;\;\square$ C $\;\;\square$ D & \multicolumn{2}{|l|}{\textbf{\instrumento}} \\
\hline
\end{tabular}
\end{center}
\singlespacing

%-------------------------------------------------------------------------
% EXPECTATIVAS DE APRENDIZAGEM 
%-------------------------------------------------------------------------

\begin{quote}
\begin{scriptsize}
\textbf{Expectativas de aprendizagem}: % (ou objetivos ou descrição dos critérios de avaliação) 
 I – Resolver corretamente e de forma completa os exercícios; II – Resolver os exercícios no tempo definido; III – Realizar a interpretação correta dos exercícios e IV – Verificar o conhecimento adquirido em sala.\\
% Outro exemplo:
\textbf{Objetivos da avaliação:}:   VERIFICAR SE O ALUNO: -Identifica as principais funções e relações trigonométricas;  identifica a aplicabilidade de teoria de seno, cosseno e tangente;  Identifica elementos do arco e opera com os mesmos;  Domina conceitos e cálculos no ciclo trigonométrico; Domina estratégias de cálculos apropriados.\\
\textbf{Instruções}:  (Exemplo: Respostas à caneta nessa folha e com as resoluções em folhas anexas. Individual, sem consulta e sem uso de calculadora ou instrumentos eletrônicos. Duração: duas aulas.)
\end{scriptsize}
\end{quote}

\hrule
\begin{center}
\begin{large}
\bf Avaliação de \disciplina \\
\end{large}
\end{center}
\hrule


%-------------------------------------------------------------------------
% QUESTÕES ou ENUNCIADOS
%-------------------------------------------------------------------------

\begin{multicols}{2}  
% Caso não queira que as questões estejam dispostas em colunas basta comentar o comando acima e o comando \end{multicols}  no final do ambiente QUESTÕES e ENUNCIADOS

\begin{enumerate} % Ambiente enumerado
\item Enunciado questão 1

\item Enunciado questão 2

\item Questão com texto matemático: O texto pode estar na linha do texto, como por exemplo: $\int x^3 ~ dx$ ou estar em linha própria: % clique sobre uma equação

$$x=\frac{-b\pm \sqrt{b^2-4ac}}{2a}$$

\item Questão com espaço em branco pré estabelecido para resposta.

\vspace{3.5cm} % Escolha o tamanho do espaço em branco mudando o valor em cm

\item Questão com figura: % É necessário fazer upload da figura

\begin{center}
\includegraphics[scale=0.2]{meBrasil.png} \\ % O nome da figura deve ser idêntico ao nome do arquivo da figura. É possível alterar a escala
{\footnotesize Legenda da Figura} % Caso não seja necessário, apague esta linha
\end{center}

\item Múltipla escolha
% Alternativas
\begin{enumerate}
    \item Alternativa 1
    \item Alternativa 2
    \item Alternativa 3
    \item Alternativa 4
    \item Alternativa 5
\end{enumerate}

\item Verdadeiro ou falso?
% Afirmações
\begin{description}
    \item[(\hspace{0.5cm})] Afirmação 1
    \item[(\hspace{0.5cm})] Afirmação 2
    \item[(\hspace{0.5cm})] Afirmação 3
    \item[(\hspace{0.5cm})] Afirmação 4
    \item[(\hspace{0.5cm})] Afirmação 5
\end{description}

\item Questão discursiva com linhas pontilhadas:
% Linhas: cada \dotfill abaixo representa uma linha pontilhada. Você pode acrescentar mais linhas ou tirá-las.

\dotfill

\dotfill

\dotfill

\dotfill

\item Questão discursiva com linhas contínuas:
%caso deseje criar linhas contínuas use \hrulefill

\hrulefill

\hrulefill

\hrulefill

\hrulefill

\end{enumerate}

\end{multicols}

%-------------------------------------------------------------------------
% MODELO DE CRITÉRIOS que podem constar na avaliação
%-------------------------------------------------------------------------
\newpage % cria a quebra de página

Leia com atenção\footnote{Confirmo que li todas as informações da atividade avaliativa com atenção e que tenho conhecimento dos critérios de avaliação descritos acima.

\

Data: \underline{\hspace{0.7cm}} / \underline{\hspace{0.7cm}} / \underline{\hspace{0.7cm}} $\quad$ Nome: \underline{\hspace{11cm}}}:

\begin{center}
    \begin{tabular}{ | l | p{15cm} |}
\multicolumn{2}{l}{}\\    
\hline
Conceito A &  Quando 	a aprendizagem foi PLENA. Nesta avaliação significa que o aluno respondeu corretamente todas as questões de forma completa e correta.\\
\hline
Conceito B & Quando 	a aprendizagem foi PARCIALMENTE PLENA. Nesta avaliação significa que o aluno respondeu correta e completamente 3 (três) questões e 1 (uma) questão parcialmente.\\
\hline
Conceito C     &  Quando 	a aprendizagem foi SUFICIENTE. Nesta avaliação significa que o aluno respondeu completa e corretamente 3 (três) questões.\\
\hline
Conceito D     & Quando 	a aprendizagem foi INSUFICIENTE. Nesta avaliação significa que o aluno NÃO 	respondeu completa e corretamente no mínimo 3(três), não atendendo de forma suficiente as expectativas de aprendizagem. \\
    \hline
    \end{tabular}
\end{center}

\end{document}